% =============================================================================
% MANUAL PENGGUNA KOMPREHENSIF
% Sistem Pengurusan Pengadil Bola Sepak Pahang (RefPahang)
% Persatuan Bola Sepak Negeri Pahang (PBNP)
% Versi 3.0 | 5 April 2026
% =============================================================================
\documentclass[12pt,a4paper]{report}

% --- Packages ---
\usepackage[utf8]{inputenc}
\usepackage[T1]{fontenc}
\usepackage[bahasam,english]{babel}
\usepackage{geometry}
\usepackage{graphicx}
\usepackage{xcolor}
\usepackage{titlesec}
\usepackage{fancyhdr}
\usepackage{hyperref}
\usepackage{tabularx}
\usepackage{longtable}
\usepackage{booktabs}
\usepackage{enumitem}
\usepackage{tcolorbox}
\usepackage{tikz}
\usepackage{fontawesome5}
\usepackage{pifont}
\usepackage{array}
\usepackage{colortbl}
\usepackage{multicol}
\usepackage{titletoc}
\usepackage{setspace}
\usepackage{parskip}
\usepackage{float}

\usetikzlibrary{arrows.meta,shapes.geometric,positioning,calc,fit}
\tcbuselibrary{skins,breakable}

% --- Page Geometry ---
\geometry{
  top=2.5cm,
  bottom=2.5cm,
  left=2.5cm,
  right=2.5cm,
  headheight=15pt
}

% --- Colors ---
\definecolor{pbnpblue}{HTML}{1E3A5F}
\definecolor{pbnpgold}{HTML}{D4A843}
\definecolor{primaryblue}{HTML}{2563EB}
\definecolor{successgreen}{HTML}{16A34A}
\definecolor{warningyellow}{HTML}{CA8A04}
\definecolor{dangered}{HTML}{DC2626}
\definecolor{infoblue}{HTML}{0284C7}
\definecolor{lightbg}{HTML}{F8FAFC}
\definecolor{tablebg}{HTML}{EFF6FF}
\definecolor{tableheader}{HTML}{1E40AF}
\definecolor{sectionblue}{HTML}{1E3A5F}
\definecolor{lightgray}{HTML}{F1F5F9}

% --- Hyperref ---
\hypersetup{
  colorlinks=true,
  linkcolor=pbnpblue,
  urlcolor=primaryblue,
  citecolor=pbnpblue,
  bookmarks=true,
  bookmarksopen=true,
  pdftitle={Manual Pengguna RefPahang v3.0},
  pdfauthor={PBNP},
  pdfsubject={Sistem Pengurusan Pengadil Bola Sepak Pahang}
}

% --- Title Formatting ---
\titleformat{\chapter}[display]
  {\normalfont\huge\bfseries\color{pbnpblue}}
  {\chaptertitlename\ \thechapter}{20pt}{\Huge}
\titlespacing*{\chapter}{0pt}{-20pt}{30pt}

\titleformat{\section}
  {\normalfont\Large\bfseries\color{sectionblue}}
  {\thesection}{1em}{}
  [\vspace{-0.5em}\textcolor{pbnpgold}{\rule{\textwidth}{1pt}}]

\titleformat{\subsection}
  {\normalfont\large\bfseries\color{pbnpblue!80}}
  {\thesubsection}{1em}{}

\titleformat{\subsubsection}
  {\normalfont\normalsize\bfseries\color{pbnpblue!70}}
  {\thesubsubsection}{1em}{}

% --- Header/Footer ---
\pagestyle{fancy}
\fancyhf{}
\fancyhead[L]{\small\textcolor{pbnpblue}{\textbf{RefPahang}} \textcolor{pbnpblue!50}{| Manual Pengguna v3.0}}
\fancyhead[R]{\small\textcolor{pbnpblue!50}{\leftmark}}
\fancyfoot[C]{\textcolor{pbnpblue}{\thepage}}
\fancyfoot[R]{\small\textcolor{pbnpblue!40}{PBNP \the\year}}
\renewcommand{\headrulewidth}{0.5pt}
\renewcommand{\headrule}{\hbox to\headwidth{\color{pbnpgold}\leaders\hrule height \headrulewidth\hfill}}
\renewcommand{\footrulewidth}{0.3pt}
\renewcommand{\footrule}{\hbox to\headwidth{\color{pbnpblue!20}\leaders\hrule height \footrulewidth\hfill}}

% --- Custom Box Environments ---
\newtcolorbox{importantbox}{
  colback=dangered!5,
  colframe=dangered!80,
  coltitle=white,
  fonttitle=\bfseries,
  title={\faExclamationTriangle\ Penting},
  rounded corners,
  boxrule=0.8pt,
  left=8pt, right=8pt, top=6pt, bottom=6pt,
  breakable
}

\newtcolorbox{infobox}{
  colback=infoblue!5,
  colframe=infoblue!60,
  coltitle=white,
  fonttitle=\bfseries,
  title={\faInfoCircle\ Nota},
  rounded corners,
  boxrule=0.8pt,
  left=8pt, right=8pt, top=6pt, bottom=6pt,
  breakable
}

\newtcolorbox{tipbox}{
  colback=successgreen!5,
  colframe=successgreen!60,
  coltitle=white,
  fonttitle=\bfseries,
  title={\faLightbulb\ Tip},
  rounded corners,
  boxrule=0.8pt,
  left=8pt, right=8pt, top=6pt, bottom=6pt,
  breakable
}

\newtcolorbox{stepbox}[1][]{
  colback=lightbg,
  colframe=pbnpblue!40,
  fonttitle=\bfseries\color{pbnpblue},
  title={#1},
  rounded corners,
  boxrule=0.5pt,
  left=8pt, right=8pt, top=6pt, bottom=6pt,
  breakable
}

% --- Custom Table Style ---
\newcolumntype{L}[1]{>{\raggedright\arraybackslash}p{#1}}
\newcolumntype{C}[1]{>{\centering\arraybackslash}p{#1}}
\newcommand{\tablemark}{\textcolor{successgreen}{\ding{51}}}
\newcommand{\tablecross}{\textcolor{dangered}{\ding{55}}}

% --- Custom Commands ---
\newcommand{\menupath}[1]{\textcolor{pbnpblue}{\textbf{#1}}}
\newcommand{\btn}[1]{\colorbox{primaryblue!10}{\textcolor{primaryblue}{\textbf{#1}}}}
\newcommand{\status}[2]{\colorbox{#1!15}{\textcolor{#1}{\small\textbf{#2}}}}
\newcommand{\rolelabel}[2]{\colorbox{#1!15}{\textcolor{#1}{\small\faUser\ \textbf{#2}}}}
\newcommand{\wajib}{\textcolor{dangered}{\small\textbf{*Wajib}}}

% --- Document ---
\begin{document}

% =============================================================================
% TITLE PAGE
% =============================================================================
\begin{titlepage}
\begin{tikzpicture}[remember picture, overlay]
  % Background
  \fill[pbnpblue] (current page.south west) rectangle (current page.north east);
  
  % Decorative elements
  \fill[pbnpblue!80] (current page.south west) -- ++(0,8) -- ++(21,0) -- ++(0,-8) -- cycle;
  \fill[pbnpgold, opacity=0.3] ([yshift=8cm]current page.south west) -- ++(0,2) -- ++(21,0) -- ++(0,-2) -- cycle;
  
  % Gold line
  \draw[pbnpgold, line width=2pt] ([yshift=10cm]current page.south west) -- ++(21,0);
  \draw[pbnpgold, line width=2pt] ([yshift=8cm]current page.south west) -- ++(21,0);
  
  % Title text
  \node[anchor=center, text=white] at ([yshift=3cm]current page.center) {
    \begin{minipage}{14cm}
      \centering
      {\fontsize{14}{16}\selectfont\textbf{PERSATUAN BOLA SEPAK NEGERI PAHANG (PBNP)}}\\[1.2cm]
      {\fontsize{36}{42}\selectfont\textbf{MANUAL PENGGUNA}}\\[0.5cm]
      {\fontsize{36}{42}\selectfont\textbf{KOMPREHENSIF}}\\[1cm]
      \textcolor{pbnpgold}{\rule{8cm}{2pt}}\\[1cm]
      {\fontsize{18}{22}\selectfont\textbf{Sistem Pengurusan Pengadil}}\\[0.3cm]
      {\fontsize{18}{22}\selectfont\textbf{Bola Sepak Pahang}}\\[0.3cm]
      {\fontsize{24}{28}\selectfont\textcolor{pbnpgold}{\textbf{RefPahang}}}\\[1.5cm]
      {\fontsize{14}{16}\selectfont Versi 3.0}
    \end{minipage}
  };
  
  % Bottom section
  \node[anchor=center, text=white] at ([yshift=3.5cm]current page.south) {
    \begin{minipage}{14cm}
      \centering
      {\fontsize{11}{13}\selectfont Tarikh Kemaskini: 5 April 2026}\\[0.3cm]
      {\fontsize{10}{12}\selectfont\textcolor{white!70}{Dokumen Rasmi --- Untuk Kegunaan Dalaman}}
    \end{minipage}
  };
\end{tikzpicture}
\end{titlepage}

% =============================================================================
% TABLE OF CONTENTS
% =============================================================================
\thispagestyle{empty}
\tableofcontents
\newpage

% =============================================================================
% BAB 1: PENGENALAN SISTEM
% =============================================================================
\chapter{Pengenalan Sistem}

\section{Apa Itu RefPahang?}

RefPahang adalah sistem pengurusan pengadil bola sepak dalam talian yang dibangunkan khusus untuk Persatuan Bola Sepak Negeri Pahang (PBNP). Sistem ini menguruskan keseluruhan kitaran hayat pengadil --- dari pendaftaran, permohonan tahunan, lantikan tugasan, rekod perlawanan, sehingga penilaian prestasi.

\section{Peranan Pengguna}

Sistem ini mempunyai \textbf{4 peranan} utama:

\begin{table}[H]
\centering
\rowcolors{2}{tablebg}{white}
\begin{tabularx}{\textwidth}{>{\bfseries\color{pbnpblue}}l l X}
\toprule
\rowcolor{tableheader}
\textcolor{white}{\textbf{Peranan}} & \textcolor{white}{\textbf{Singkatan}} & \textcolor{white}{\textbf{Tanggungjawab Utama}} \\
\midrule
Pentadbir & Admin & Pengurusan keseluruhan sistem, kelulusan permohonan, lantikan pengadil, tetapan \\
PP Daerah & PP & Pegawai Pembangunan Daerah --- pengesahan permohonan \& perlawanan daerah masing-masing \\
Pengadil & --- & Pengadil bola sepak --- pendaftaran, permohonan tahunan, rekod perlawanan, terima lantikan \\
Penilai Pengadil & RA & Referee Assessor --- menilai prestasi pengadil semasa perlawanan \\
\bottomrule
\end{tabularx}
\caption{Peranan pengguna dalam sistem RefPahang}
\end{table}

\section{Akses Sistem}

\begin{table}[H]
\centering
\rowcolors{2}{tablebg}{white}
\begin{tabularx}{\textwidth}{l X}
\toprule
\rowcolor{tableheader}
\textcolor{white}{\textbf{Maklumat}} & \textcolor{white}{\textbf{Nilai}} \\
\midrule
URL Sistem & \texttt{https://refpahang.com} \\
Pelayar Disokong & Google Chrome, Mozilla Firefox, Safari, Microsoft Edge (versi terkini) \\
Peranti & Desktop, tablet, dan telefon pintar (responsif) \\
\bottomrule
\end{tabularx}
\end{table}

\section{Saluran Notifikasi}

Sistem menghantar notifikasi melalui \textbf{3 saluran}:

\begin{enumerate}[label=\textcolor{pbnpblue}{\textbf{\arabic*.}}, leftmargin=2em]
  \item \textbf{Portal} --- Notifikasi dalam sistem (ikon loceng di bahagian atas)
  \item \textbf{Emel} --- Dihantar ke emel berdaftar pengguna
  \item \textbf{Telegram} --- Dihantar melalui bot Telegram RefPahang (perlu dihubungkan terlebih dahulu)
\end{enumerate}

% =============================================================================
% BAB 2: PANDUAN UMUM
% =============================================================================
\chapter{Panduan Umum (Semua Pengguna)}

\section{Pendaftaran Akaun Baharu}

\subsection{Langkah-langkah}

\begin{stepbox}[Langkah 1: Buka Halaman Pendaftaran]
Buka laman \texttt{https://refpahang.com} $\rightarrow$ klik \btn{Daftar Akaun Baharu}
\end{stepbox}

\begin{stepbox}[Langkah 2: Isi Borang Pendaftaran]
\begin{table}[H]
\centering
\rowcolors{2}{tablebg}{white}
\begin{tabularx}{\textwidth}{l X c}
\toprule
\rowcolor{tableheader}
\textcolor{white}{\textbf{Medan}} & \textcolor{white}{\textbf{Penerangan}} & \textcolor{white}{\textbf{Wajib}} \\
\midrule
Nama Penuh & Seperti dalam kad pengenalan & \tablemark \\
No. Kad Pengenalan & 12 digit tanpa sengkang (cth: 900101065432) & \tablemark \\
Emel & Emel aktif untuk menerima notifikasi & \tablemark \\
No. Telefon & 10--11 digit (cth: 0123456789) & \tablemark \\
Jantina & Lelaki / Perempuan & \tablemark \\
Jenis Pengadil & Pengadil Kebangsaan / Negeri / Penilai / PP & \tablemark \\
Persatuan Bolasepak Daerah & Pilih daerah (12 daerah di Pahang) & \tablemark \\
Pengesahan Data & Tandakan kotak persetujuan & \tablemark \\
\bottomrule
\end{tabularx}
\end{table}
\end{stepbox}

\begin{stepbox}[Langkah 3--5: Penghantaran]
\begin{enumerate}[label=\textcolor{pbnpblue}{\textbf{\arabic*.}}, start=3]
  \item Klik \btn{Daftar}
  \item \textbf{Kata laluan} akan dijana secara automatik dan dihantar ke emel anda
  \item Modal kejayaan akan dipaparkan --- klik \btn{Buka Telegram} untuk menghubungkan akaun Telegram
\end{enumerate}
\end{stepbox}

\begin{importantbox}
Sila tukar kata laluan anda selepas log masuk pertama kali.
\end{importantbox}

\subsection{12 Persatuan Bolasepak Daerah (PBD)}

\begin{table}[H]
\centering
\rowcolors{2}{tablebg}{white}
\begin{tabularx}{\textwidth}{c l c l}
\toprule
\rowcolor{tableheader}
\textcolor{white}{\textbf{\#}} & \textcolor{white}{\textbf{Daerah}} & \textcolor{white}{\textbf{\#}} & \textcolor{white}{\textbf{Daerah}} \\
\midrule
1 & Bentong & 7 & Maran \\
2 & Bera & 8 & Pekan \\
3 & Cameron Highlands & 9 & Raub \\
4 & Jerantut & 10 & Rompin \\
5 & Kuantan & 11 & Temerloh \\
6 & Lipis & 12 & Muadzam Shah \\
\bottomrule
\end{tabularx}
\caption{Senarai 12 PBD Pahang}
\end{table}

\subsection{Penentuan Peranan}

\begin{itemize}[label=\textcolor{pbnpblue}{\faChevronRight}]
  \item Pilih \textbf{``Pegawai Pembangunan''} $\rightarrow$ Akaun didaftarkan sebagai \rolelabel{primaryblue}{PP Daerah}
  \item Pilih jenis lain $\rightarrow$ Akaun didaftarkan sebagai \rolelabel{successgreen}{Pengadil}
\end{itemize}

\begin{infobox}
Pendaftaran Penilai Pengadil (RA) dilakukan melalui jenis ``Penilai Pengadil'' tetapi peranan akaun tetap Pengadil. Admin akan menaikkan taraf ke peranan Penilai jika perlu.
\end{infobox}

\section{Log Masuk}

\begin{enumerate}[label=\textcolor{pbnpblue}{\textbf{Langkah \arabic*:}}, leftmargin=3.5em]
  \item Buka \texttt{https://refpahang.com} $\rightarrow$ halaman log masuk
  \item Masukkan \textbf{Emel} dan \textbf{Kata Laluan}
  \item Klik \btn{Log Masuk}
  \item Sistem mengarahkan anda ke dashboard mengikut peranan:
    \begin{itemize}
      \item Admin $\rightarrow$ \texttt{/admin}
      \item PP Daerah $\rightarrow$ \texttt{/pp-daerah}
      \item Pengadil $\rightarrow$ \texttt{/pengadil}
      \item Penilai $\rightarrow$ \texttt{/penilai}
    \end{itemize}
\end{enumerate}

\begin{infobox}
Jika ini adalah log masuk pertama, anda akan diminta menukar kata laluan.
\end{infobox}

\section{Lupa Kata Laluan}

\begin{enumerate}[label=\textcolor{pbnpblue}{\textbf{\arabic*.}}, leftmargin=2em]
  \item Di halaman log masuk, klik \btn{Lupa kata laluan?}
  \item Masukkan emel berdaftar anda
  \item Klik \btn{Hantar}
  \item Pautan penetapan semula akan dihantar ke emel anda
  \item Klik pautan dalam emel $\rightarrow$ tetapkan kata laluan baharu
\end{enumerate}

\begin{infobox}
\textbf{Nota Keselamatan:} Mesej kejayaan sentiasa dipaparkan walaupun emel tidak wujud dalam sistem --- ini untuk melindungi privasi pengguna.
\end{infobox}

\section{Tukar Kata Laluan}

\begin{enumerate}[label=\textcolor{pbnpblue}{\textbf{\arabic*.}}, leftmargin=2em]
  \item Pergi ke \menupath{Profil Saya} $\rightarrow$ bahagian \textbf{Tukar Kata Laluan}
  \item Masukkan:
    \begin{itemize}
      \item \textbf{Kata Laluan Semasa} --- kata laluan yang sedang digunakan
      \item \textbf{Kata Laluan Baharu} --- minimum 8 aksara
      \item \textbf{Sahkan Kata Laluan Baharu} --- mestilah sama dengan di atas
    \end{itemize}
  \item Klik \btn{Tukar Kata Laluan}
\end{enumerate}

\section{Menghubungkan Telegram}

Telegram membolehkan anda menerima notifikasi segera untuk lantikan, permohonan, dan maklumat penting.

\begin{enumerate}[label=\textcolor{pbnpblue}{\textbf{\arabic*.}}, leftmargin=2em]
  \item Pergi ke \menupath{Profil Saya}
  \item Cari bahagian \textbf{Telegram} $\rightarrow$ klik \btn{Hubungkan Telegram}
  \item Sistem akan membuka aplikasi Telegram dengan bot RefPahang
  \item Tekan \btn{Start} atau hantar mesej \texttt{/start} kepada bot
  \item Bot akan mengesahkan penyambungan --- Status bertukar menjadi \status{successgreen}{Telegram Berjaya Dihubungkan}
\end{enumerate}

\begin{tipbox}
Anda juga boleh menghubungkan Telegram sejurus selepas pendaftaran melalui modal kejayaan.
\end{tipbox}

\section{Notifikasi}

\subsection{Melihat Notifikasi}
\begin{enumerate}[label=\textcolor{pbnpblue}{\textbf{\arabic*.}}, leftmargin=2em]
  \item Klik ikon \textbf{loceng} \faBell\ di bahagian atas
  \item Senarai 50 notifikasi terkini dipaparkan
  \item Bilangan notifikasi belum dibaca ditunjukkan pada lencana merah
\end{enumerate}

\subsection{Menanda Dibaca}
\begin{itemize}[label=\textcolor{pbnpblue}{\faChevronRight}]
  \item Klik pada notifikasi individu $\rightarrow$ ditanda dibaca
  \item Klik \btn{Tanda Semua Dibaca} $\rightarrow$ semua notifikasi ditanda
\end{itemize}

\subsection{Jenis Notifikasi}

\begin{table}[H]
\centering
\rowcolors{2}{tablebg}{white}
\begin{tabularx}{\textwidth}{l X}
\toprule
\rowcolor{tableheader}
\textcolor{white}{\textbf{Jenis}} & \textcolor{white}{\textbf{Penerangan}} \\
\midrule
Pendaftaran Berjaya & Akaun baharu berjaya didaftarkan \\
Lantikan Baru & Tugasan lantikan baharu diterima \\
Permohonan Disahkan & Permohonan disahkan oleh PP / Admin \\
Permohonan Ditolak & Permohonan ditolak \\
Perlawanan Baru & Rekod perlawanan baharu perlu disahkan (untuk PP) \\
Pengesahan Perlawanan & Perlawanan disahkan/ditolak oleh PP \\
Rekod Perlawanan & Anda direkodkan dalam perlawanan oleh pengadil lain \\
Profil Dikemaskini & Profil berjaya dikemaskini \\
\bottomrule
\end{tabularx}
\caption{Jenis notifikasi dalam sistem}
\end{table}

\section{Kemaskini Profil}

Semua pengguna boleh mengemaskini maklumat profil mereka.

\subsection{Maklumat Yang Boleh Dikemaskini}

\textbf{Maklumat Peribadi:}
\begin{itemize}[label=\textcolor{pbnpblue}{\faChevronRight}, nosep]
  \item Nama Penuh, No. Telefon, Saiz Baju, Tahun Mula Aktif
  \item Alamat Penuh (Alamat 1, Alamat 2, Poskod, Daerah)
\end{itemize}

\textbf{Maklumat Pekerjaan:}
\begin{itemize}[label=\textcolor{pbnpblue}{\faChevronRight}, nosep]
  \item Status Kerja, Jawatan, Nama Majikan
  \item Alamat Majikan (Alamat 1, Alamat 2, Poskod, Daerah, Negeri)
\end{itemize}

\textbf{Maklumat Waris:}
\begin{itemize}[label=\textcolor{pbnpblue}{\faChevronRight}, nosep]
  \item Nama Waris, Hubungan (Ibu Bapa/Pasangan/Anak/Adik-beradik/Lain-lain), No. Telefon Waris
\end{itemize}

\textbf{Gambar Profil:}
\begin{itemize}[label=\textcolor{pbnpblue}{\faChevronRight}, nosep]
  \item Klik ikon kamera $\rightarrow$ pilih gambar $\rightarrow$ potong (crop) $\rightarrow$ simpan
  \item Format: JPEG, PNG, WebP | Saiz maksimum: 5MB
\end{itemize}

\begin{importantbox}
\textbf{Maklumat yang TIDAK boleh diubah:} No. Kad Pengenalan, Emel, Jantina, Negeri (ditetapkan Pahang)
\end{importantbox}

% =============================================================================
% BAB 3: MANUAL PENGADIL
% =============================================================================
\chapter{Manual Pengadil}

\section{Menu Utama Pengadil}

\begin{table}[H]
\centering
\rowcolors{2}{tablebg}{white}
\begin{tabularx}{\textwidth}{c l c X}
\toprule
\rowcolor{tableheader}
\textcolor{white}{\textbf{\#}} & \textcolor{white}{\textbf{Menu}} & \textcolor{white}{\textbf{Ikon}} & \textcolor{white}{\textbf{Fungsi}} \\
\midrule
1 & Dashboard & \faChartBar & Paparan ringkasan dan statistik \\
2 & Profil Saya & \faUser & Pengurusan maklumat peribadi \\
3 & Tugasan Lantikan & \faClipboardList & Senarai lantikan dari Admin \\
4 & Rekod Perlawanan & \faFutbol & Rekod perlawanan tidak rasmi \\
5 & Penilaian Saya & \faStar & Laporan penilaian prestasi \\
6 & \textbf{Permohonan} & \faEdit & --- \\
6a & \quad Pendaftaran Tahunan & \faTag & Borang pendaftaran tahunan \\
6b & \quad Ujian Kecergasan & \faDumbbell & Maklumat ujian kecergasan \\
6c & \quad Ujian Kelas III FAM & \faEdit & Permohonan ujian Kelas III \\
6d & \quad Ujian Kelas I FAM & \faMedal & Permohonan ujian Kelas I \\
\bottomrule
\end{tabularx}
\caption{Menu utama Pengadil}
\end{table}

\section{Dashboard Pengadil}

Dashboard memaparkan ringkasan maklumat penting:

\subsection{Kad Statistik (4 kad)}

\begin{table}[H]
\centering
\rowcolors{2}{tablebg}{white}
\begin{tabularx}{\textwidth}{l X}
\toprule
\rowcolor{tableheader}
\textcolor{white}{\textbf{Kad}} & \textcolor{white}{\textbf{Penerangan}} \\
\midrule
\textbf{Perlawanan (Tahun)} & Jumlah perlawanan tahun semasa \\
\textbf{Disahkan PP} & Bilangan perlawanan yang telah disahkan oleh PP Daerah \\
\textbf{Status Kelayakan} & ``Layak'' ($\geq$20 perlawanan disahkan) / ``Belum Layak'' \\
\textbf{Lantikan Belum Jawab} & Bilangan tugasan lantikan yang belum dijawab \\
\bottomrule
\end{tabularx}
\end{table}

\subsection{Bahagian Lain}
\begin{itemize}[label=\textcolor{pbnpblue}{\faChevronRight}]
  \item \textbf{Pautan Pantas} --- Profil, Tugasan Lantikan, Rekod Perlawanan, Pengadil Berdaftar
  \item \textbf{Notifikasi Terkini} --- 5 notifikasi terbaru
  \item \textbf{Pengumuman} --- Pengumuman dari PBNP
  \item \textbf{Perlawanan Terkini} --- Jadual 5 perlawanan terakhir
\end{itemize}

\section{Permohonan Pendaftaran Tahunan}

\subsection{Tujuan}
Pengadil perlu membuat pendaftaran tahunan setiap tahun untuk kekal aktif. Pendaftaran ini mewajibkan bayaran yuran dan penyerahan dokumen.

\subsection{Syarat}
\begin{itemize}[label=\textcolor{dangered}{\faExclamationCircle}]
  \item Akaun sudah didaftarkan dalam sistem
  \item Profil telah dilengkapi (alamat, pekerjaan, waris)
  \item \textbf{Minimum 20 perlawanan disahkan} dalam tahun semasa (untuk pengadil sahaja)
  \item Tempoh pendaftaran dibuka oleh Admin
\end{itemize}

\subsection{Langkah-langkah}

\begin{stepbox}[Langkah 1: Buat Bayaran]
Buat bayaran ke akaun yang dipaparkan (nama bank, no. akaun, jumlah yuran). Simpan resit bayaran (PDF/gambar).
\end{stepbox}

\begin{stepbox}[Langkah 2: Isi Borang]
\begin{itemize}[nosep]
  \item Maklumat peribadi diisi secara automatik dari profil
  \item Isi/kemaskini \textbf{Maklumat Waris}: Nama, Hubungan, No. Telefon \wajib
  \item Pilih \textbf{Saiz Baju Rasmi}: XS, S, M, L, XL, XXL, XXXL \wajib
  \item Muat naik \textbf{Resit Bayaran}: PDF/JPG/JPEG/PNG, maks. 5MB \wajib
  \item Muat naik \textbf{Gambar Profil Terbaru}: JPG/JPEG/PNG, maks. 5MB \wajib
\end{itemize}
\end{stepbox}

\begin{stepbox}[Langkah 3: Deklarasi Kesihatan (5 item --- semua wajib)]
\begin{enumerate}[label=\textcolor{successgreen}{\faCheckSquare}\ , nosep]
  \item Status kesihatan fizikal dan mental baik
  \item Telah melalui pemeriksaan perubatan dan disahkan layak
  \item Memahami risiko dan melepaskan PBNP dari tuntutan
  \item Bersetuju memaklumkan sebarang kondisi perubatan
  \item Memberi kebenaran rawatan kecemasan
\end{enumerate}
\end{stepbox}

\begin{stepbox}[Langkah 4: Perakuan Umum (3 item --- semua wajib)]
\begin{enumerate}[label=\textcolor{successgreen}{\faCheckSquare}\ , nosep]
  \item Maklumat yang diberikan adalah benar dan tepat
  \item Sekiranya maklumat palsu, permohonan akan dibatalkan
  \item Bersetuju mematuhi peraturan dan undang-undang PBNP/FAM
\end{enumerate}
\end{stepbox}

\begin{stepbox}[Langkah 5: Hantar]
Klik \btn{Hantar Permohonan}
\end{stepbox}

\subsection{Aliran Kelulusan}

\begin{center}
\begin{tikzpicture}[
  node distance=1.2cm,
  process/.style={rectangle, rounded corners, draw=pbnpblue, fill=pbnpblue!10, text width=3.8cm, minimum height=1cm, text centered, font=\small},
  decision/.style={diamond, draw=warningyellow!80!black, fill=warningyellow!10, text width=2cm, minimum height=1cm, text centered, font=\small, aspect=2},
  success/.style={rectangle, rounded corners, draw=successgreen, fill=successgreen!10, text width=2.8cm, minimum height=0.8cm, text centered, font=\small\bfseries},
  fail/.style={rectangle, rounded corners, draw=dangered, fill=dangered!10, text width=2.2cm, minimum height=0.8cm, text centered, font=\small\bfseries},
  arrow/.style={-{Stealth[length=6pt]}, thick, pbnpblue}
]
  \node[process] (submit) {Pengadil hantar permohonan};
  \node[process, below=of submit] (waitpp) {Status: Menunggu PP Daerah};
  \node[decision, below=of waitpp] (pp) {PP Daerah};
  \node[process, below left=1.5cm and 0.5cm of pp] (waitadmin) {Status: Menunggu Admin};
  \node[fail, below right=1.5cm and 0.5cm of pp] (reject1) {Ditolak};
  \node[decision, below=of waitadmin] (admin) {Admin};
  \node[success, below left=1.5cm and 0.3cm of admin] (approve) {Lengkap \faCheck};
  \node[fail, below right=1.5cm and 0.3cm of admin] (reject2) {Ditolak};
  
  \draw[arrow] (submit) -- (waitpp);
  \draw[arrow] (waitpp) -- (pp);
  \draw[arrow] (pp) -- node[left, font=\scriptsize] {Sahkan} (waitadmin);
  \draw[arrow] (pp) -- node[right, font=\scriptsize] {Tolak} (reject1);
  \draw[arrow] (waitadmin) -- (admin);
  \draw[arrow] (admin) -- node[left, font=\scriptsize] {Luluskan} (approve);
  \draw[arrow] (admin) -- node[right, font=\scriptsize] {Tolak} (reject2);
\end{tikzpicture}
\end{center}

\section{Permohonan Ujian Kelas III FAM}

\subsection{Tujuan}
Memohon untuk menduduki Ujian Kelas III FAM (ujian bertulis).

\subsection{Syarat Kelayakan}
\begin{itemize}[label=\textcolor{pbnpblue}{\faChevronRight}]
  \item Umur: \textbf{15 hingga 40 tahun} (dikira dari No. KP)
  \item Yuran: \textbf{RM 50}
  \item Tempoh permohonan dibuka oleh Admin
\end{itemize}

\subsection{Langkah-langkah}
\begin{enumerate}[label=\textcolor{pbnpblue}{\textbf{\arabic*.}}, leftmargin=2em]
  \item Buat bayaran RM 50 ke akaun FAM yang dipaparkan
  \item Isi maklumat waris
  \item Muat naik resit bayaran FAM
  \item Tandakan 2 perakuan umum
  \item Klik \btn{Hantar Permohonan}
\end{enumerate}

\section{Permohonan Ujian Kelas I FAM}

\subsection{Tujuan}
Memohon untuk menduduki Ujian Kelas I FAM (peningkatan taraf).

\subsection{Syarat Kelayakan}
\begin{itemize}[label=\textcolor{pbnpblue}{\faChevronRight}]
  \item Umur: \textbf{Tidak melebihi 32 tahun}
  \item Telah \textbf{lulus Ujian Kelas III} sekurang-kurangnya \textbf{2 tahun} yang lalu
  \item Yuran: \textbf{RM 300}
  \item Tempoh permohonan dibuka oleh Admin
\end{itemize}

\begin{importantbox}
Sistem akan menyemak secara automatik kelayakan umur dan tempoh lulus Kelas III.
\end{importantbox}

\section{Sejarah Permohonan}

Di setiap halaman permohonan, bahagian \textbf{``Sejarah Permohonan''} memaparkan:

\begin{table}[H]
\centering
\rowcolors{2}{tablebg}{white}
\begin{tabularx}{\textwidth}{l X}
\toprule
\rowcolor{tableheader}
\textcolor{white}{\textbf{Lajur}} & \textcolor{white}{\textbf{Penerangan}} \\
\midrule
Tarikh Hantar & Tarikh permohonan dihantar \\
Sesi & Tahun permohonan \\
Status & Menunggu / Diluluskan / Ditolak / Lulus / Tidak Lulus / Tidak Hadir \\
Catatan & Nota dari PP Daerah atau Admin \\
Muat Turun & Pautan muat turun borang PDF (untuk permohonan yang diluluskan) \\
\bottomrule
\end{tabularx}
\end{table}

\section{Tugasan Lantikan}

\subsection{Tujuan}
Menerima dan menjawab tugasan lantikan rasmi dari Admin untuk kejohanan/perlawanan.

\subsection{Maklumat Dipaparkan}

\begin{table}[H]
\centering
\rowcolors{2}{tablebg}{white}
\begin{tabularx}{\textwidth}{l X}
\toprule
\rowcolor{tableheader}
\textcolor{white}{\textbf{Medan}} & \textcolor{white}{\textbf{Penerangan}} \\
\midrule
Kejohanan & Nama kejohanan \\
No. Perlawanan & Nombor perlawanan \\
Tarikh \& Masa & Tarikh dan masa perlawanan \\
Pasukan & Pasukan tuan rumah vs tetamu \\
Tempat & Lokasi perlawanan \\
Jawatan & Pengadil / Penolong Pengadil 1 / Penolong Pengadil 2 / Pegawai Ke-4 \\
Pengadil Utama & Nama pengadil utama (untuk rujukan) \\
Status & Belum Jawab / Diterima / Ditolak \\
\bottomrule
\end{tabularx}
\end{table}

\subsection{Menjawab Tugasan}

\subsubsection{Melalui Portal (dalam sistem)}
\begin{enumerate}[label=\textcolor{pbnpblue}{\textbf{\arabic*.}}, leftmargin=2em]
  \item Pergi ke \menupath{Tugasan Lantikan}
  \item Cari tugasan dengan status \status{warningyellow}{Belum Jawab}
  \item Klik \btn{Terima} atau \btn{Tolak}
    \begin{itemize}
      \item Jika tolak: Wajib isi sebab penolakan
    \end{itemize}
  \item Status akan dikemaskini serta-merta
\end{enumerate}

\subsubsection{Melalui Pautan Emel/Telegram (tanpa log masuk)}
\begin{enumerate}[label=\textcolor{pbnpblue}{\textbf{\arabic*.}}, leftmargin=2em]
  \item Terima notifikasi melalui emel atau Telegram
  \item Klik pautan \btn{Terima} atau \btn{Tolak} dalam mesej
  \item Halaman pengesahan akan dipaparkan
  \item Pautan adalah \textbf{sekali guna} --- selepas digunakan, ia tamat tempoh
\end{enumerate}

\subsection{Apa Yang Berlaku Selepas Terima}
\begin{itemize}[label=\textcolor{successgreen}{\faCheckCircle}]
  \item Rekod perlawanan \textbf{dicipta secara automatik} dengan status \status{successgreen}{Disahkan}
  \item Jika anda seorang Penilai Pengadil, \textbf{token borang penilaian} akan dijana dan dihantar
  \item Jika \textbf{semua pegawai} untuk perlawanan tersebut telah menerima, status jadual perlawanan bertukar kepada \status{successgreen}{Disahkan}
\end{itemize}

\section{Rekod Perlawanan}

\subsection{Tujuan}
Merekod perlawanan \textbf{tidak rasmi} (persahabatan, perlawanan luar kejohanan) sebagai minit perlawanan. Perlawanan rasmi dari lantikan dicipta secara automatik dan tidak perlu didaftar di sini.

\subsection{Menambah Perlawanan Baharu}

\begin{importantbox}
Hanya \textbf{seorang pengadil} perlu mengisi borang untuk \textbf{SEMUA pegawai} dalam perlawanan tersebut. Rekod akan dipaparkan kepada semua pegawai yang terlibat.
\end{importantbox}

\begin{stepbox}[Langkah 1: Maklumat Perlawanan]
\begin{table}[H]
\centering
\rowcolors{2}{tablebg}{white}
\begin{tabularx}{\textwidth}{l l c X}
\toprule
\rowcolor{tableheader}
\textcolor{white}{\textbf{Medan}} & \textcolor{white}{\textbf{Jenis}} & \textcolor{white}{\textbf{Wajib}} & \textcolor{white}{\textbf{Pilihan}} \\
\midrule
Pasukan Tuan Rumah & Teks & \tablemark & --- \\
Pasukan Lawan & Teks & \tablemark & --- \\
Tarikh & Tarikh & \tablemark & --- \\
Masa & Masa & \tablecross & --- \\
Jenis & Pilihan & \tablecross & Liga, Piala, Persahabatan, Lain \\
Nama Kejohanan & Teks & \tablecross & cth: Liga M-League 2026 \\
Tempat & Teks & \tablecross & Lokasi perlawanan \\
Daerah Perlawanan & Pilihan & \tablemark & 12 daerah Pahang \\
Cuaca & Pilihan & \tablecross & Cerah, Mendung, Hujan Renyai, dll. \\
\bottomrule
\end{tabularx}
\end{table}
\end{stepbox}

\begin{stepbox}[Langkah 2: Keputusan Perlawanan]
\begin{table}[H]
\centering
\rowcolors{2}{tablebg}{white}
\begin{tabularx}{\textwidth}{l X}
\toprule
\rowcolor{tableheader}
\textcolor{white}{\textbf{Tempoh}} & \textcolor{white}{\textbf{Penerangan}} \\
\midrule
Separuh Masa (HT) & Skor home : away \\
Penuh Masa (FT) & Skor home : away \\
Masa Tambahan (ET) & Skor home : away (opsional) \\
Penalti (PS) & Skor home : away (opsional) \\
\bottomrule
\end{tabularx}
\end{table}
\end{stepbox}

\begin{stepbox}[Langkah 3: Pegawai Perlawanan]
\begin{itemize}[nosep]
  \item Sistem menyediakan \textbf{4 slot} lalai: Pengadil, Penolong Pengadil 1, Penolong Pengadil 2, Pengadil Ke-4
  \item Boleh tambah hingga \textbf{5 pegawai} (termasuk Pengadil Ke-5)
  \item Untuk setiap pegawai: Pilih \textbf{Jawatan}, kemudian \textbf{cari pengadil} (nama atau No. KP)
  \item Klik \btn{+ Tambah Pegawai} untuk menambah slot
\end{itemize}
\end{stepbox}

\subsection{Peraturan Penting}

\begin{table}[H]
\centering
\rowcolors{2}{tablebg}{white}
\begin{tabularx}{\textwidth}{>{\bfseries}l X}
\toprule
\rowcolor{tableheader}
\textcolor{white}{\textbf{Peraturan}} & \textcolor{white}{\textbf{Penerangan}} \\
\midrule
Peraturan 14 Hari & Perlawanan tidak boleh didaftarkan jika tarikh melebihi 14 hari yang lalu \\
Jawatan Unik & Setiap jawatan hanya boleh diberikan kepada seorang pegawai \\
Minimum 1 Pegawai & Sekurang-kurangnya seorang pegawai perlu diisi \\
Daerah Wajib & Daerah perlawanan wajib dipilih --- PP Daerah tersebut akan menerima notifikasi \\
Tidak Boleh Padam Lantikan & Rekod yang dicipta dari lantikan rasmi tidak boleh dipadamkan \\
\bottomrule
\end{tabularx}
\end{table}

\section{Penilaian Saya}

\subsection{Tujuan}
Melihat laporan penilaian prestasi yang telah disahkan oleh Admin.

\subsection{Pengekodan Warna Markah}

\begin{table}[H]
\centering
\begin{tabularx}{\textwidth}{l l X}
\toprule
\rowcolor{tableheader}
\textcolor{white}{\textbf{Julat}} & \textcolor{white}{\textbf{Warna}} & \textcolor{white}{\textbf{Penerangan}} \\
\midrule
\rowcolor{successgreen!10} 8.3 -- 10.0 & \textcolor{successgreen}{\faCircle\ Hijau} & Prestasi baik hingga sangat baik \\
\rowcolor{primaryblue!10} 8.0 -- 8.2 & \textcolor{primaryblue}{\faCircle\ Biru} & Baik dengan penambahbaikan \\
\rowcolor{warningyellow!10} 7.5 -- 7.9 & \textcolor{warningyellow}{\faCircle\ Kuning} & Tidak memuaskan \\
\rowcolor{dangered!10} $<$ 7.5 & \textcolor{dangered}{\faCircle\ Merah} & Buruk \\
\bottomrule
\end{tabularx}
\caption{Pengekodan warna markah penilaian}
\end{table}

\subsection{Butiran Laporan}
\begin{itemize}[label=\textcolor{pbnpblue}{\faChevronRight}]
  \item Maklumat perlawanan penuh (kejohanan, pasukan, skor, tarikh, penilai)
  \item Jadual semua pegawai: Jawatan, Nama, Markah, Prestasi
  \item \textbf{Penilaian anda:} Markah, prestasi, \textcolor{successgreen}{Kekuatan (tag hijau)}, \textcolor{dangered}{Kelemahan (tag merah)}, Nasihat
  \item Ulasan Keseluruhan dan Catatan Admin
  \item Butang \btn{Muat Turun PDF}
\end{itemize}

% =============================================================================
% BAB 4: MANUAL PP DAERAH
% =============================================================================
\chapter{Manual PP Daerah}

\section{Menu Utama PP Daerah}

\begin{table}[H]
\centering
\rowcolors{2}{tablebg}{white}
\begin{tabularx}{\textwidth}{c l X}
\toprule
\rowcolor{tableheader}
\textcolor{white}{\textbf{\#}} & \textcolor{white}{\textbf{Menu}} & \textcolor{white}{\textbf{Fungsi}} \\
\midrule
1 & Dashboard & Ringkasan daerah \\
2 & \textbf{Permohonan Saya} & --- \\
2a & \quad Pendaftaran Tahunan (R4) & Permohonan sendiri \\
2b & \quad Ujian Kelas III FAM & Permohonan sendiri \\
2c & \quad Ujian Kelas I FAM & Permohonan sendiri \\
3 & Pengesahan Permohonan & Sahkan/tolak permohonan daerah \\
4 & Pengesahan Perlawanan & Sahkan/tolak perlawanan daerah \\
5 & \textbf{Permohonan} & --- \\
5a & \quad Pendaftaran Tahunan (R4) & Senarai permohonan daerah \\
5b & \quad Ujian Kelas III FAM & Senarai permohonan daerah \\
5c & \quad Ujian Kelas I FAM & Senarai permohonan daerah \\
6 & Pengadil Berdaftar & Senarai pengadil daerah \\
7 & RA Berdaftar & Senarai penilai daerah \\
8 & Laporan Penilaian & Laporan penilaian daerah \\
9 & Tugasan Lantikan & Tugasan lantikan sendiri \\
10 & Profil & Pengurusan profil sendiri \\
\bottomrule
\end{tabularx}
\caption{Menu utama PP Daerah}
\end{table}

\begin{importantbox}
\textbf{Skop Penting:} Semua data yang dipaparkan kepada PP Daerah \textbf{hanya melibatkan daerah sendiri}. PP Kuantan hanya nampak data Kuantan.
\end{importantbox}

\section{Dashboard PP Daerah}

\subsection{Kad Statistik (4 kad)}
\begin{table}[H]
\centering
\rowcolors{2}{tablebg}{white}
\begin{tabularx}{\textwidth}{l X}
\toprule
\rowcolor{tableheader}
\textcolor{white}{\textbf{Kad}} & \textcolor{white}{\textbf{Penerangan}} \\
\midrule
Pengadil Berdaftar & Jumlah pengadil aktif dalam daerah \\
Menunggu Pengesahan & Permohonan yang perlu disahkan (animasi denyut jika $>$ 0) \\
Diluluskan & Permohonan yang telah diluluskan \\
Perlawanan Bulan Ini & Jumlah perlawanan bulan semasa \\
\bottomrule
\end{tabularx}
\end{table}

\subsection{Bahagian Lain}
\begin{itemize}[label=\textcolor{pbnpblue}{\faChevronRight}]
  \item \textbf{Sepanduk ambar} jika ada permohonan menunggu $\rightarrow$ pautan ke Pengesahan Permohonan
  \item \textbf{Sepanduk biru} jika ada perlawanan menunggu $\rightarrow$ pautan ke Pengesahan Perlawanan
  \item \textbf{Top 5 Pengadil} --- Pengadil paling aktif (berdasarkan jumlah perlawanan)
  \item \textbf{Perlawanan Akan Datang} --- 10 perlawanan seterusnya dengan pegawai yang dilantik
\end{itemize}

\section{Pengesahan Permohonan}

\subsection{Tujuan}
PP Daerah berperanan sebagai pengesah peringkat pertama sebelum permohonan dihantar ke Admin.

\subsection{Aliran Proses}

\begin{center}
\begin{tikzpicture}[
  node distance=1.2cm,
  process/.style={rectangle, rounded corners, draw=pbnpblue, fill=pbnpblue!10, text width=4cm, minimum height=0.9cm, text centered, font=\small},
  arrow/.style={-{Stealth[length=6pt]}, thick, pbnpblue}
]
  \node[process] (wait) {Menunggu PP Daerah};
  \node[process, below left=1.5cm and 1cm of wait] (sahkan) {Menunggu Admin};
  \node[process, below right=1.5cm and 1cm of wait] (tolak1) {\textcolor{dangered}{Ditolak}};
  \node[process, below left=1.5cm and 0.2cm of sahkan] (lulus) {\textcolor{successgreen}{Lengkap \faCheck}};
  \node[process, below right=1.5cm and 0.2cm of sahkan] (tolak2) {\textcolor{dangered}{Ditolak}};
  
  \draw[arrow] (wait) -- node[left, font=\scriptsize] {PP Sahkan} (sahkan);
  \draw[arrow] (wait) -- node[right, font=\scriptsize] {PP Tolak} (tolak1);
  \draw[arrow] (sahkan) -- node[left, font=\scriptsize] {Admin Luluskan} (lulus);
  \draw[arrow] (sahkan) -- node[right, font=\scriptsize] {Admin Tolak} (tolak2);
\end{tikzpicture}
\end{center}

\subsection{Mengesahkan Permohonan}
\begin{enumerate}[label=\textcolor{pbnpblue}{\textbf{\arabic*.}}, leftmargin=2em]
  \item Semak maklumat pengadil dan resit bayaran
  \item Klik \btn{Sahkan} $\rightarrow$ Dialog pengesahan $\rightarrow$ Klik \btn{Ya}
  \item Permohonan bertukar kepada status \status{primaryblue}{Menunggu Admin}
  \item \textbf{Emel dihantar kepada} semua Admin aktif dan pemohon
\end{enumerate}

\subsection{Menolak Permohonan}
\begin{enumerate}[label=\textcolor{pbnpblue}{\textbf{\arabic*.}}, leftmargin=2em]
  \item Klik \btn{Tolak} $\rightarrow$ Modal penolakan dipaparkan
  \item \textbf{Wajib} isi sebab penolakan dalam kotak teks
  \item Klik \btn{Hantar Penolakan}
  \item Permohonan bertukar kepada status \status{dangered}{Ditolak}
  \item \textbf{Emel dihantar kepada pemohon} --- memberitahu beserta sebab
\end{enumerate}

\section{Pengesahan Perlawanan}

\subsection{Tujuan}
PP Daerah mengesahkan atau menolak rekod perlawanan tidak rasmi yang didaftarkan oleh pengadil.

\begin{infobox}
\textbf{Skop Pengesahan:} Perlawanan berkumpulan disenaraikan berdasarkan \textbf{daerah perlawanan} yang dipilih oleh pengadil (bukan daerah pengadil). Jika pengadil dari Kuantan mendaftar perlawanan di Pekan, PP Pekan yang perlu mengesahkan.
\end{infobox}

\subsection{Tindakan PP}

\begin{table}[H]
\centering
\rowcolors{2}{tablebg}{white}
\begin{tabularx}{\textwidth}{l l X}
\toprule
\rowcolor{tableheader}
\textcolor{white}{\textbf{Tindakan}} & \textcolor{white}{\textbf{Ikon}} & \textcolor{white}{\textbf{Penerangan}} \\
\midrule
Sahkan & \textcolor{successgreen}{\faCheck} & Semua rekod kumpulan disahkan sekaligus \\
Tolak & \textcolor{dangered}{\faTimes} & Semua rekod kumpulan ditolak \\
Edit Semula & \textcolor{primaryblue}{\faRedo} & Kembalikan ke ``Belum Disahkan'' \\
Padam & \textcolor{dangered}{\faTrash} & Padam semua rekod kumpulan (tidak boleh batal) \\
\bottomrule
\end{tabularx}
\end{table}

\section{Senarai Pengadil Berdaftar}

Melihat semua pengadil berdaftar dalam daerah. Klik pada nama untuk melihat profil penuh:

\begin{itemize}[label=\textcolor{pbnpblue}{\faChevronRight}, nosep]
  \item Maklumat Peribadi, Alamat, Maklumat Pengadil
  \item Pekerjaan, Waris, Statistik
  \item Rekod Perlawanan Terkini (10 perlawanan terakhir)
\end{itemize}

\section{Laporan Penilaian}

\subsection{Statistik}
4 angka ringkasan: Jumlah Laporan, Purata Markah, Markah Tertinggi, Markah Terendah

\subsection{Butiran Laporan}
\begin{itemize}[label=\textcolor{pbnpblue}{\faChevronRight}]
  \item Maklumat perlawanan penuh dan keputusan (HT, FT, ET, PS)
  \item Jadual semua pegawai: Jawatan, Nama, Markah, Prestasi
  \item Penilaian per pegawai: Kekuatan, Kelemahan, Nasihat (per bahagian)
  \item Ulasan Keseluruhan
  \item Butang \btn{Muat Turun PDF}
\end{itemize}

% =============================================================================
% BAB 5: MANUAL PENILAI PENGADIL (RA)
% =============================================================================
\chapter{Manual Penilai Pengadil (RA)}

\section{Menu Utama Penilai}

\begin{table}[H]
\centering
\rowcolors{2}{tablebg}{white}
\begin{tabularx}{\textwidth}{c l X}
\toprule
\rowcolor{tableheader}
\textcolor{white}{\textbf{\#}} & \textcolor{white}{\textbf{Menu}} & \textcolor{white}{\textbf{Fungsi}} \\
\midrule
1 & Dashboard & Ringkasan tugasan dan statistik \\
2 & Penilaian Pengadil & Isi dan hantar laporan penilaian \\
3 & Pendaftaran Tahunan (R4) & Permohonan pendaftaran tahunan \\
4 & Tugasan & Senarai tugasan lantikan \\
5 & Statistik & Statistik penilaian \\
6 & Profil & Pengurusan profil sendiri \\
\bottomrule
\end{tabularx}
\caption{Menu utama Penilai}
\end{table}

\section{Penilaian Pengadil}

\subsection{Tujuan}
Ini adalah fungsi utama Penilai --- menilai prestasi pengadil semasa perlawanan kejohanan.

\subsection{Status Laporan}

\begin{table}[H]
\centering
\begin{tabularx}{\textwidth}{l l X}
\toprule
\rowcolor{tableheader}
\textcolor{white}{\textbf{Status}} & \textcolor{white}{\textbf{Warna}} & \textcolor{white}{\textbf{Penerangan}} \\
\midrule
\rowcolor{dangered!8} Belum Dilengkapi & \textcolor{dangered}{\faCircle} & Belum mula mengisi \\
\rowcolor{warningyellow!8} Draf & \textcolor{warningyellow}{\faCircle} & Sudah mula, belum dihantar \\
\rowcolor{primaryblue!8} Dihantar & \textcolor{primaryblue}{\faCircle} & Sudah dihantar ke Admin \\
\rowcolor{successgreen!8} Disahkan & \textcolor{successgreen}{\faCircle} & Admin telah mengesahkan \\
\bottomrule
\end{tabularx}
\caption{Status laporan penilaian}
\end{table}

\subsection{Mengisi Borang Penilaian}

\subsubsection{Kaedah Akses}

\begin{table}[H]
\centering
\rowcolors{2}{tablebg}{white}
\begin{tabularx}{\textwidth}{l X}
\toprule
\rowcolor{tableheader}
\textcolor{white}{\textbf{Kaedah}} & \textcolor{white}{\textbf{Penerangan}} \\
\midrule
Melalui Portal & Log masuk $\rightarrow$ Penilaian Pengadil $\rightarrow$ Perlu Dinilai $\rightarrow$ Mula Penilaian \\
Melalui Pautan Token & Klik pautan dalam emel/Telegram --- tanpa perlu log masuk \\
\bottomrule
\end{tabularx}
\end{table}

\subsection{Bahagian Borang Penilaian}

\subsubsection{1. Maklumat Perlawanan (automatik)}
Kejohanan, Pasukan, No. Perlawanan, Tarikh, Masa, Tempat, Nama Penilai

\subsubsection{2. Keputusan Perlawanan}
Separuh Masa (HT), Penuh Masa (FT), Masa Tambahan (ET), Penalti (PS)

\subsubsection{3. Tahap Kesukaran Perlawanan}
Normal / Susah / Sangat Susah

\subsubsection{4. Cuaca}
Cerah / Mendung / Hujan Renyai / Hujan Lebat / Panas Terik / Berangin

\subsubsection{5. Penilaian Per Pegawai}

\textbf{Kriteria Berbeza Mengikut Jawatan:}

\begin{table}[H]
\centering
\rowcolors{2}{tablebg}{white}
\begin{tabularx}{\textwidth}{l l c X}
\toprule
\rowcolor{tableheader}
\textcolor{white}{\textbf{Jawatan}} & \textcolor{white}{\textbf{Bahagian}} & \textcolor{white}{\textbf{Bil.}} & \textcolor{white}{\textbf{Contoh}} \\
\midrule
\multirow{3}{*}{\textbf{Pengadil (R)}} & Kawalan Permainan & 28 & Penalti, kad kuning/merah, bola tangan \\
 & Kecergasan Fizikal & 10 & Kelajuan larian, stamina, penempatan \\
 & Kerjasama Berpasukan & 12 & Komunikasi, konsistensi keputusan \\
\midrule
\textbf{Penolong Pengadil} & Penilaian Penolong & 24 & Keputusan offside, isyarat, dsb. \\
\midrule
\textbf{Pegawai Ke-4} & Penilaian P4 & 11 & Pengurusan teknikal, dsb. \\
\bottomrule
\end{tabularx}
\caption{Kriteria penilaian mengikut jawatan}
\end{table}

Untuk setiap bahagian:
\begin{itemize}[label=\textcolor{pbnpblue}{\faChevronRight}, nosep]
  \item \textcolor{successgreen}{\textbf{Kekuatan}} (pilih pelbagai) --- kriteria yang pengadil lakukan dengan baik
  \item \textcolor{dangered}{\textbf{Kelemahan}} (pilih pelbagai) --- kriteria yang perlu diperbaiki
  \item \textbf{Nasihat} (teks) --- saranan khusus untuk penambahbaikan
\end{itemize}

\subsection{Skala Pemarkahan (Rujukan)}

\begin{table}[H]
\centering
\begin{tabularx}{\textwidth}{l X}
\toprule
\rowcolor{tableheader}
\textcolor{white}{\textbf{Julat Markah}} & \textcolor{white}{\textbf{Penerangan}} \\
\midrule
\rowcolor{successgreen!10} 9.0 -- 10.0 & Prestasi sangat baik, keputusan pada situasi sukar adalah betul \\
\rowcolor{successgreen!8} 8.5 -- 8.9 & Prestasi sangat baik \\
\rowcolor{successgreen!5} 8.3 -- 8.4 & Prestasi baik \\
\rowcolor{primaryblue!8} 8.0 -- 8.2 & Baik dengan penambahbaikan \\
\rowcolor{warningyellow!8} 7.9 & Tidak memuaskan --- 1 insiden penting (sepatutnya 8.3+) \\
\rowcolor{warningyellow!10} 7.8 & Tidak memuaskan --- 1 insiden penting (sepatutnya 8.0--8.2) \\
\rowcolor{warningyellow!12} 7.5 -- 7.7 & Tidak memuaskan --- 2 insiden penting \\
\rowcolor{dangered!8} 7.4 & Kesilapan mempengaruhi keputusan perlawanan \\
\rowcolor{dangered!12} 7.0 -- 7.4 & Buruk --- 3 atau lebih insiden penting \\
\bottomrule
\end{tabularx}
\caption{Skala pemarkahan penilaian pengadil}
\end{table}

\subsection{Simpan dan Hantar}

\begin{table}[H]
\centering
\rowcolors{2}{tablebg}{white}
\begin{tabularx}{\textwidth}{l X}
\toprule
\rowcolor{tableheader}
\textcolor{white}{\textbf{Tindakan}} & \textcolor{white}{\textbf{Penerangan}} \\
\midrule
\textbf{Simpan Draf} & Markah dan maklumat disimpan tetapi belum dihantar. Boleh sambung kemudian. \\
\textbf{Hantar Laporan} & Semua markah pegawai \textbf{WAJIB} diisi. Status $\rightarrow$ ``Dihantar'' $\rightarrow$ menunggu Admin. \\
\bottomrule
\end{tabularx}
\end{table}

\subsection{Selepas Admin Sahkan}
\begin{itemize}[label=\textcolor{successgreen}{\faCheckCircle}]
  \item Status bertukar kepada \status{successgreen}{Disahkan}
  \item \textbf{Emel diterima} oleh setiap pegawai yang dinilai (markah, 5 kekuatan, 5 kelemahan, pautan PDF)
  \item \textbf{Mesej Telegram} dihantar kepada pegawai yang telah menghubungkan Telegram
  \item Laporan boleh dimuat turun sebagai PDF
\end{itemize}

% =============================================================================
% BAB 6: MANUAL PENTADBIR (ADMIN)
% =============================================================================
\chapter{Manual Pentadbir (Admin)}

\section{Menu Utama Admin}

\begin{table}[H]
\centering
\rowcolors{2}{tablebg}{white}
\begin{tabularx}{\textwidth}{c l X}
\toprule
\rowcolor{tableheader}
\textcolor{white}{\textbf{\#}} & \textcolor{white}{\textbf{Menu}} & \textcolor{white}{\textbf{Fungsi}} \\
\midrule
1 & Dashboard & Paparan ringkasan keseluruhan sistem \\
2 & Permohonan Pengadil & Urus permohonan semua pengadil \\
3 & Permohonan RA & Urus permohonan penilai \\
4 & Pengadil Berdaftar & Senarai semua pengadil berdaftar \\
5 & RA Berdaftar & Senarai semua penilai berdaftar \\
6 & PP Daerah & Senarai semua PP Daerah \\
7 & Pengguna & Pengurusan semua pengguna \\
8 & Lantikan Pengadil & Kejohanan, jadual \& lantikan \\
9 & Pengadil Luar & Urus pengadil dari negeri lain \\
10 & Statistik & Analitik dan laporan \\
11 & Pengumuman & Urus pengumuman awam \\
12 & Tetapan & Tetapan sistem \\
13 & Profil Saya & Pengurusan profil Admin \\
\bottomrule
\end{tabularx}
\caption{Menu utama Admin}
\end{table}

\section{Dashboard Admin}

\subsection{Kad KPI (4 kad)}
\begin{table}[H]
\centering
\rowcolors{2}{tablebg}{white}
\begin{tabularx}{\textwidth}{l X}
\toprule
\rowcolor{tableheader}
\textcolor{white}{\textbf{Kad}} & \textcolor{white}{\textbf{Penerangan}} \\
\midrule
Pengadil Aktif & Jumlah pengadil aktif dalam sistem \\
Permohonan Baru & Permohonan menunggu tindakan (animasi denyut jika $>$ 0) \\
Laporan Semakan & Laporan penilaian yang belum disahkan \\
Perlawanan Bulan Ini & Jumlah perlawanan dalam bulan semasa \\
\bottomrule
\end{tabularx}
\end{table}

\subsection{Widget dan Akses Pantas}
\begin{itemize}[label=\textcolor{pbnpblue}{\faChevronRight}]
  \item \textbf{Perhatian Segera:} Permohonan tertunda + Laporan perlu disemak
  \item \textbf{Permohonan Terkini:} 5 terbaru dengan butang Luluskan/Tolak terus
  \item \textbf{Akses Pantas:} Pengumuman, Statistik
\end{itemize}

\section{Permohonan Pengadil}

\subsection{Tab Permohonan}
\begin{table}[H]
\centering
\rowcolors{2}{tablebg}{white}
\begin{tabularx}{\textwidth}{l X}
\toprule
\rowcolor{tableheader}
\textcolor{white}{\textbf{Tab}} & \textcolor{white}{\textbf{Fungsi}} \\
\midrule
Pengadil Berdaftar & Pendaftaran tahunan pengadil \\
Ujian Kelas III FAM & Permohonan ujian bertulis Kelas III \\
Ujian Kelas I FAM & Permohonan ujian Kelas I \\
\bottomrule
\end{tabularx}
\end{table}

\subsection{Tindakan Permohonan}

\begin{table}[H]
\centering
\rowcolors{2}{tablebg}{white}
\begin{tabularx}{\textwidth}{l X}
\toprule
\rowcolor{tableheader}
\textcolor{white}{\textbf{Tindakan}} & \textcolor{white}{\textbf{Penerangan}} \\
\midrule
\textcolor{successgreen}{\textbf{Luluskan}} & Status $\rightarrow$ ``Lengkap''. Emel pengesahan dihantar kepada pemohon. \\
\textcolor{dangered}{\textbf{Tolak}} & Wajib isi sebab. Status $\rightarrow$ ``Ditolak''. Emel penolakan dihantar. \\
\textcolor{primaryblue}{\textbf{Lulus}} (Ujian) & Calon lulus ujian \\
\textcolor{dangered}{\textbf{Tidak Lulus}} (Ujian) & Calon tidak lulus ujian \\
\textcolor{warningyellow}{\textbf{Tidak Hadir}} (Ujian) & Calon tidak hadir ujian \\
\bottomrule
\end{tabularx}
\end{table}

\section{Pengadil Berdaftar / RA Berdaftar}

\begin{itemize}[label=\textcolor{pbnpblue}{\faChevronRight}]
  \item \textbf{Lihat Profil} --- Modal butiran penuh (3 tab: Maklumat, Sejarah, Perlawanan)
  \item \textbf{Edit} --- Kemaskini maklumat pengadil
  \item \textbf{Padam} --- Padam rekod (dengan pengesahan)
  \item \textbf{Muat Turun Excel} --- Eksport senarai
  \item \textbf{Cetak} --- Susun atur mesra cetakan
\end{itemize}

\section{Pengurusan Pengguna}

\begin{table}[H]
\centering
\rowcolors{2}{tablebg}{white}
\begin{tabularx}{\textwidth}{l X}
\toprule
\rowcolor{tableheader}
\textcolor{white}{\textbf{Tindakan}} & \textcolor{white}{\textbf{Penerangan}} \\
\midrule
Tambah Pengguna & Cipta akaun baharu \\
Edit Profil & Kemaskini maklumat pengguna \\
Reset Kata Laluan & Hantar pautan reset \\
Hantar Semula Emel & Hantar semula emel notifikasi \\
Padam Pengguna & Padam akaun pengguna \\
\bottomrule
\end{tabularx}
\end{table}

\section{Lantikan Pengadil (5 Tab)}

Ini adalah modul paling kompleks dalam sistem --- mengurus keseluruhan aliran lantikan dari kejohanan hingga penilaian.

\subsection{Tab 1: Kejohanan}

Mendaftar dan mengurus kejohanan bola sepak.

\begin{itemize}[label=\textcolor{pbnpblue}{\faChevronRight}]
  \item \textbf{Tambah Kejohanan} --- nama, tarikh mula/akhir, anjuran, logo
  \item \textbf{Edit / Padam} --- urus kejohanan
  \item \textbf{Lihat Jadual} --- pergi ke Tab 2
\end{itemize}

\subsection{Tab 2: Jadual Perlawanan}

Mengurus jadual perlawanan dalam kejohanan.

\begin{itemize}[label=\textcolor{pbnpblue}{\faChevronRight}]
  \item \textbf{Tambah Manual} --- isi butiran per perlawanan
  \item \textbf{Muat Naik Excel} --- muat naik jadual secara pukal
  \item \textbf{Edit / Padam} --- urus perlawanan
\end{itemize}

\subsection{Tab 3: Lantikan Pengadil}

\begin{stepbox}[Aliran Kerja Lantikan]
\begin{enumerate}[label=\textcolor{pbnpblue}{\textbf{\arabic*.}}, leftmargin=2em]
  \item \textbf{Pilih perlawanan} dari senarai jadual
  \item Sistem memaparkan \textbf{slot jawatan}: Pengadil, Penolong Pengadil 1 \& 2, Pegawai Ke-4, Penilai Pengadil (RA)
  \item Lihat \textbf{Pool Pengadil} --- senarai pengadil layak (berdaftar + luar)
  \item \textbf{Klik pengadil} $\rightarrow$ pilih slot $\rightarrow$ Modal pengesahan
  \item Klik \btn{Lantik} $\rightarrow$ Rekod lantikan dicipta
  \item Pengadil menerima notifikasi: \textbf{Portal + Emel + Telegram}
\end{enumerate}
\end{stepbox}

\subsection{Tab 4: Jadual Lantikan}

Laporan penuh lantikan untuk keseluruhan kejohanan.

\begin{itemize}[label=\textcolor{pbnpblue}{\faChevronRight}]
  \item Semua perlawanan dengan pegawai dan status respons
  \item \textbf{Pengesahan Jadual}: Isi nama pengesah, jawatan, nota
  \item \textbf{Muat Turun PDF} --- laporan A4
\end{itemize}

\subsection{Tab 5: Laporan Penilaian}

Melihat dan mengesahkan laporan penilaian daripada Penilai.

\begin{itemize}[label=\textcolor{pbnpblue}{\faChevronRight}]
  \item \textbf{Sahkan Laporan} $\rightarrow$ Emel \& Telegram dihantar kepada pegawai
  \item \textbf{Muat Turun PDF} --- laporan A4
\end{itemize}

\section{Statistik \& Analitik}

\subsection{Tab Permohonan}
\begin{itemize}[label=\textcolor{pbnpblue}{\faChevronRight}, nosep]
  \item Carta Bar: Permohonan mengikut jenis borang
  \item Carta Pai: Taburan status permohonan
  \item Carta Garisan: Trend bulanan permohonan
\end{itemize}

\subsection{Tab Pengadil}
\begin{itemize}[label=\textcolor{pbnpblue}{\faChevronRight}, nosep]
  \item Carta Donat: Taburan jantina
  \item Carta Bar: Status pekerjaan \& taburan daerah
  \item Carta Garisan: Trend pendaftaran bulanan
\end{itemize}

\subsection{Tab Lantikan}
\begin{itemize}[label=\textcolor{pbnpblue}{\faChevronRight}, nosep]
  \item Kad: Jumlah, Selesai, Disahkan, Belum
  \item Jadual: Senarai kejohanan dengan statistik
\end{itemize}

\section{Pengumuman}

\begin{table}[H]
\centering
\rowcolors{2}{tablebg}{white}
\begin{tabularx}{\textwidth}{l X}
\toprule
\rowcolor{tableheader}
\textcolor{white}{\textbf{Tindakan}} & \textcolor{white}{\textbf{Penerangan}} \\
\midrule
Tambah & Modal borang: tajuk + kandungan \\
Lihat & Modal paparan penuh \\
Edit & Kemaskini tajuk / kandungan \\
Padam & Padam dengan pengesahan \\
\bottomrule
\end{tabularx}
\end{table}

\begin{infobox}
Pengumuman dipaparkan di halaman log masuk (carousel), dashboard semua peranan, dan bahagian pengumuman portal.
\end{infobox}

\section{Tetapan Sistem}

\subsection{Mod Penyelenggaraan}
\begin{itemize}[label=\textcolor{pbnpblue}{\faChevronRight}]
  \item \textbf{Aktif}: Hanya Admin boleh akses sistem
  \item \textbf{Tidak Aktif}: Semua pengguna boleh akses
  \item Togol suis untuk tukar
\end{itemize}

\subsection{Status Permohonan Per Jenis}

\begin{table}[H]
\centering
\rowcolors{2}{tablebg}{white}
\begin{tabularx}{\textwidth}{l X}
\toprule
\rowcolor{tableheader}
\textcolor{white}{\textbf{Jenis}} & \textcolor{white}{\textbf{Kawalan}} \\
\midrule
Daftar Akaun Baru & DIBUKA / DITUTUP + tarikh mula \& akhir \\
Pendaftaran Tahunan (R1+R2+R4) & DIBUKA / DITUTUP + tarikh mula \& akhir \\
Ujian Kelas III FAM & DIBUKA / DITUTUP + tarikh mula \& akhir \\
Ujian Kelas I FAM & DIBUKA / DITUTUP + tarikh mula \& akhir \\
Pendaftaran Penilai (R4) & DIBUKA / DITUTUP + tarikh mula \& akhir \\
\bottomrule
\end{tabularx}
\end{table}

\subsection{Maklumat Pembayaran}
Nama Bank, No. Akaun, Jumlah Yuran

\subsection{Peraturan Kelayakan FAM}
Umur minimum/maksimum ujian, Bilangan pusingan kecergasan

% =============================================================================
% BAB 7: CARTA ALIR PROSES
% =============================================================================
\chapter{Carta Alir Proses}

\section{Aliran Pendaftaran Akaun Baharu}

\begin{center}
\begin{tikzpicture}[
  node distance=1cm,
  process/.style={rectangle, rounded corners=4pt, draw=pbnpblue, fill=pbnpblue!8, text width=9cm, minimum height=0.9cm, text centered, font=\small},
  decision/.style={diamond, draw=warningyellow!80!black, fill=warningyellow!8, text width=3cm, minimum height=0.8cm, text centered, font=\small, aspect=2.5},
  arrow/.style={-{Stealth[length=5pt]}, thick, pbnpblue!70}
]
  \node[process] (a) {Pengguna baru buka halaman pendaftaran};
  \node[process, below=of a] (b) {Isi borang (Nama, IC, Emel, Telefon, Jantina, Jenis, PBD)};
  \node[process, below=of b] (c) {Klik ``Daftar''};
  \node[decision, below=1.2cm of c] (d) {Jenis?};
  \node[process, below left=1.2cm and 0.3cm of d, text width=4cm] (e1) {Peranan: PP Daerah};
  \node[process, below right=1.2cm and 0.3cm of d, text width=4cm] (e2) {Peranan: Pengadil};
  \node[process, below=3cm of d] (f) {Kata laluan dijana automatik $\rightarrow$ dihantar ke emel};
  \node[process, below=of f] (g) {Modal kejayaan $\rightarrow$ Pilihan hubungkan Telegram};
  \node[process, below=of g] (h) {Log masuk $\rightarrow$ Tukar kata laluan $\rightarrow$ Lengkapkan profil};
  
  \draw[arrow] (a) -- (b);
  \draw[arrow] (b) -- (c);
  \draw[arrow] (c) -- (d);
  \draw[arrow] (d) -| node[above left, font=\scriptsize] {Pegawai Pembangunan} (e1);
  \draw[arrow] (d) -| node[above right, font=\scriptsize] {Jenis lain} (e2);
  \draw[arrow] (e1) |- (f);
  \draw[arrow] (e2) |- (f);
  \draw[arrow] (f) -- (g);
  \draw[arrow] (g) -- (h);
\end{tikzpicture}
\end{center}

\section{Aliran Permohonan Pendaftaran Tahunan}

\begin{center}
\begin{tikzpicture}[
  node distance=1cm,
  process/.style={rectangle, rounded corners=4pt, draw=pbnpblue, fill=pbnpblue!8, text width=8cm, minimum height=0.9cm, text centered, font=\small},
  decision/.style={diamond, draw=warningyellow!80!black, fill=warningyellow!8, text width=2.2cm, minimum height=0.8cm, text centered, font=\small, aspect=2.5},
  success/.style={rectangle, rounded corners=4pt, draw=successgreen, fill=successgreen!10, text width=3cm, minimum height=0.8cm, text centered, font=\small\bfseries},
  fail/.style={rectangle, rounded corners=4pt, draw=dangered, fill=dangered!10, text width=2.5cm, minimum height=0.8cm, text centered, font=\small\bfseries},
  arrow/.style={-{Stealth[length=5pt]}, thick, pbnpblue!70}
]
  \node[process] (a) {Pastikan profil lengkap + $\geq$20 perlawanan disahkan};
  \node[process, below=of a] (b) {Buat bayaran yuran tahunan};
  \node[process, below=of b] (c) {Isi borang + Deklarasi + Hantar permohonan};
  \node[process, below=of c] (d) {Status: ``Menunggu PP Daerah''};
  \node[decision, below=1.2cm of d] (pp) {PP Daerah};
  \node[process, below left=1.5cm and 0.5cm of pp, text width=4.5cm] (waitadmin) {Status: ``Menunggu Admin''};
  \node[fail, below right=1.5cm and 0.5cm of pp] (rej1) {Ditolak \faTimes};
  \node[decision, below=1.2cm of waitadmin] (admin) {Admin};
  \node[success, below left=1.2cm and 0cm of admin] (ok) {Lengkap \faCheck};
  \node[fail, below right=1.2cm and 0cm of admin] (rej2) {Ditolak \faTimes};
  
  \draw[arrow] (a) -- (b);
  \draw[arrow] (b) -- (c);
  \draw[arrow] (c) -- (d);
  \draw[arrow] (d) -- (pp);
  \draw[arrow] (pp) -- node[left, font=\scriptsize] {Sahkan} (waitadmin);
  \draw[arrow] (pp) -- node[right, font=\scriptsize] {Tolak} (rej1);
  \draw[arrow] (waitadmin) -- (admin);
  \draw[arrow] (admin) -- node[left, font=\scriptsize] {Luluskan} (ok);
  \draw[arrow] (admin) -- node[right, font=\scriptsize] {Tolak} (rej2);
\end{tikzpicture}
\end{center}

\section{Aliran Lantikan Pengadil}

\begin{center}
\begin{tikzpicture}[
  node distance=1cm,
  process/.style={rectangle, rounded corners=4pt, draw=pbnpblue, fill=pbnpblue!8, text width=9cm, minimum height=0.9cm, text centered, font=\small},
  decision/.style={diamond, draw=warningyellow!80!black, fill=warningyellow!8, text width=2.5cm, minimum height=0.8cm, text centered, font=\small, aspect=2.5},
  success/.style={rectangle, rounded corners=4pt, draw=successgreen, fill=successgreen!10, text width=7cm, minimum height=0.8cm, text centered, font=\small},
  fail/.style={rectangle, rounded corners=4pt, draw=dangered, fill=dangered!10, text width=5cm, minimum height=0.8cm, text centered, font=\small},
  arrow/.style={-{Stealth[length=5pt]}, thick, pbnpblue!70}
]
  \node[process] (a) {Admin cipta Kejohanan $\rightarrow$ Muat Naik Jadual Perlawanan};
  \node[process, below=of a] (b) {Pilih perlawanan $\rightarrow$ Lantik pengadil dari pool ke slot jawatan};
  \node[process, below=of b] (c) {Notifikasi dihantar: Portal + Emel + Telegram};
  \node[decision, below=1.2cm of c] (d) {Pengadil};
  \node[success, below left=1.5cm and -1cm of d] (e1) {TERIMA $\rightarrow$ Rekod perlawanan dicipta automatik};
  \node[fail, below right=1.5cm and -0.5cm of d] (e2) {TOLAK (wajib isi sebab)};
  \node[success, below=of e1] (f) {Token penilaian dijana untuk Penilai (jika berkenaan)};
  \node[success, below=of f] (g) {Jika SEMUA terima $\rightarrow$ Jadual perlawanan: ``Disahkan''};
  
  \draw[arrow] (a) -- (b);
  \draw[arrow] (b) -- (c);
  \draw[arrow] (c) -- (d);
  \draw[arrow] (d) -- node[left, font=\scriptsize] {Terima} (e1);
  \draw[arrow] (d) -- node[right, font=\scriptsize] {Tolak} (e2);
  \draw[arrow] (e1) -- (f);
  \draw[arrow] (f) -- (g);
\end{tikzpicture}
\end{center}

\section{Aliran Rekod Perlawanan (Tidak Rasmi)}

\begin{center}
\begin{tikzpicture}[
  node distance=1cm,
  process/.style={rectangle, rounded corners=4pt, draw=pbnpblue, fill=pbnpblue!8, text width=9cm, minimum height=0.9cm, text centered, font=\small},
  decision/.style={diamond, draw=warningyellow!80!black, fill=warningyellow!8, text width=2cm, minimum height=0.8cm, text centered, font=\small, aspect=2.5},
  success/.style={rectangle, rounded corners=4pt, draw=successgreen, fill=successgreen!10, text width=4cm, minimum height=0.8cm, text centered, font=\small\bfseries},
  fail/.style={rectangle, rounded corners=4pt, draw=dangered, fill=dangered!10, text width=4cm, minimum height=0.8cm, text centered, font=\small\bfseries},
  arrow/.style={-{Stealth[length=5pt]}, thick, pbnpblue!70}
]
  \node[process] (a) {Pengadil klik ``Tambah Perlawanan''};
  \node[process, below=of a] (b) {Isi: Pasukan, Tarikh, Jenis, Tempat, Daerah, Cuaca, Keputusan};
  \node[process, below=of b] (c) {Pilih pegawai perlawanan (1--5 orang) + jawatan};
  \node[process, below=of c] (d) {Klik ``Simpan'' $\rightarrow$ Rekod dicipta untuk SETIAP pegawai};
  \node[process, below=of d] (e) {PP Daerah (daerah perlawanan) menerima notifikasi};
  \node[decision, below=1.2cm of e] (f) {PP Daerah};
  \node[success, below left=1.2cm and 0cm of f] (g1) {Disahkan \faCheck};
  \node[fail, below right=1.2cm and 0cm of f] (g2) {Tidak Disahkan \faTimes};
  
  \draw[arrow] (a) -- (b);
  \draw[arrow] (b) -- (c);
  \draw[arrow] (c) -- (d);
  \draw[arrow] (d) -- (e);
  \draw[arrow] (e) -- (f);
  \draw[arrow] (f) -- node[left, font=\scriptsize] {Sahkan} (g1);
  \draw[arrow] (f) -- node[right, font=\scriptsize] {Tolak} (g2);
\end{tikzpicture}
\end{center}

\section{Aliran Penilaian Pengadil}

\begin{center}
\begin{tikzpicture}[
  node distance=1cm,
  process/.style={rectangle, rounded corners=4pt, draw=pbnpblue, fill=pbnpblue!8, text width=9cm, minimum height=0.9cm, text centered, font=\small},
  success/.style={rectangle, rounded corners=4pt, draw=successgreen, fill=successgreen!10, text width=9cm, minimum height=0.8cm, text centered, font=\small},
  arrow/.style={-{Stealth[length=5pt]}, thick, pbnpblue!70}
]
  \node[process] (a) {Admin melantik Penilai Pengadil (RA) ke perlawanan};
  \node[process, below=of a] (b) {Penilai terima lantikan $\rightarrow$ Token penilaian dijana};
  \node[process, below=of b] (c) {Penilai mengisi borang (portal atau pautan token)};
  \node[process, below=of c] (d) {Isi: Keputusan, Kesukaran, Cuaca, Markah per pegawai};
  \node[process, below=of d] (e) {``Simpan Draf'' atau ``Hantar Laporan''};
  \node[process, below=of e] (f) {Admin menyemak laporan $\rightarrow$ ``Sahkan''};
  \node[success, below=of f] (g) {Emel + Telegram kepada pengadil (markah, PDF)};
  \node[success, below=of g] (h) {Pengadil melihat markah di ``Penilaian Saya''};
  
  \draw[arrow] (a) -- (b);
  \draw[arrow] (b) -- (c);
  \draw[arrow] (c) -- (d);
  \draw[arrow] (d) -- (e);
  \draw[arrow] (e) -- (f);
  \draw[arrow] (f) -- (g);
  \draw[arrow] (g) -- (h);
\end{tikzpicture}
\end{center}

% =============================================================================
% BAB 8: PENYELESAIAN MASALAH (TROUBLESHOOT)
% =============================================================================
\chapter{Penyelesaian Masalah (Troubleshoot)}

\section{Masalah Log Masuk}

\begin{longtable}{L{4cm} L{3.5cm} L{6.5cm}}
\toprule
\rowcolor{tableheader}
\textcolor{white}{\textbf{Masalah}} & \textcolor{white}{\textbf{Punca}} & \textcolor{white}{\textbf{Penyelesaian}} \\
\midrule
\endhead
\rowcolor{tablebg}
``Emel atau kata laluan tidak sah'' & Kata laluan salah atau emel tidak berdaftar & Semak ejaan emel. Guna fungsi ``Lupa Kata Laluan'' untuk tetapkan semula \\
Tidak boleh log masuk langsung & Akaun tidak aktif & Hubungi Admin untuk mengaktifkan semula akaun \\
\rowcolor{tablebg}
Halaman ``Sistem Dalam Penyelenggaraan'' & Mod penyelenggaraan aktif & Hanya Admin boleh akses. Tunggu sehingga mod dimatikan \\
Selepas log masuk, kembali ke halaman log masuk & Sesi tamat atau cookies disekat & Pastikan cookies dibenarkan. Cuba pelayar lain. Buka dalam mod biasa \\
\bottomrule
\end{longtable}

\section{Masalah Pendaftaran}

\begin{longtable}{L{4cm} L{3.5cm} L{6.5cm}}
\toprule
\rowcolor{tableheader}
\textcolor{white}{\textbf{Masalah}} & \textcolor{white}{\textbf{Punca}} & \textcolor{white}{\textbf{Penyelesaian}} \\
\midrule
\endhead
\rowcolor{tablebg}
``Pendaftaran Ditutup'' & Admin telah menutup pendaftaran & Hubungi Admin atau tunggu dibuka \\
``No. KP sudah digunakan'' & Akaun dengan IC sama sudah wujud & Guna ``Lupa Kata Laluan'' untuk akses \\
\rowcolor{tablebg}
``Emel sudah digunakan'' & Emel sudah berdaftar & Cuba emel lain atau guna ``Lupa Kata Laluan'' \\
Tidak terima emel kata laluan & Emel dalam folder spam & Semak folder spam. Hubungi Admin jika perlu \\
\bottomrule
\end{longtable}

\section{Masalah Permohonan}

\begin{longtable}{L{4cm} L{3.5cm} L{6.5cm}}
\toprule
\rowcolor{tableheader}
\textcolor{white}{\textbf{Masalah}} & \textcolor{white}{\textbf{Punca}} & \textcolor{white}{\textbf{Penyelesaian}} \\
\midrule
\endhead
\rowcolor{tablebg}
``Pendaftaran ditutup'' & Tempoh sudah tamat & Hubungi Admin untuk tarikh pembukaan \\
``Minimum 20 perlawanan'' & Belum cukup disahkan & Pastikan $\geq$20 disahkan oleh PP \\
\rowcolor{tablebg}
``Umur tidak memenuhi syarat'' & Umur di luar lingkungan & Syarat automatik dari No. KP \\
``Perlu lulus Kelas III 2 tahun'' & Prasyarat belum dipenuhi & Pastikan lulus $\geq$2 tahun lalu \\
\rowcolor{tablebg}
``Permohonan sudah dihantar'' & 1 per jenis per tahun & Tunggu permohonan semasa diproses \\
Fail tidak boleh dimuat naik & Format/saiz tidak sesuai & Guna PDF/JPG/PNG, maks 5MB \\
\rowcolor{tablebg}
Status ``Ditolak'' & PP/Admin menolak & Semak catatan, perbaiki, hantar semula \\
\bottomrule
\end{longtable}

\section{Masalah Perlawanan}

\begin{longtable}{L{4cm} L{3.5cm} L{6.5cm}}
\toprule
\rowcolor{tableheader}
\textcolor{white}{\textbf{Masalah}} & \textcolor{white}{\textbf{Punca}} & \textcolor{white}{\textbf{Penyelesaian}} \\
\midrule
\endhead
\rowcolor{tablebg}
``Rekod melebihi 14 hari'' & Tarikh $>$ 14 hari lalu & Mesti daftar dalam 14 hari \\
``Jawatan sudah digunakan'' & 2 pegawai jawatan sama & 1 jawatan = 1 pegawai \\
\rowcolor{tablebg}
Tidak jumpa pengadil & Belum berdaftar & Pastikan ada akaun aktif \\
``Sila pilih daerah'' & Daerah tidak dipilih & Pilih daerah lokasi perlawanan \\
\rowcolor{tablebg}
Tidak boleh padam & Rekod lantikan rasmi & Lantikan rasmi tidak boleh dipadam \\
``Tidak Disahkan'' & PP menolak & Semak catatan PP, hubungi PP \\
\bottomrule
\end{longtable}

\section{Masalah Lantikan}

\begin{longtable}{L{4cm} L{3.5cm} L{6.5cm}}
\toprule
\rowcolor{tableheader}
\textcolor{white}{\textbf{Masalah}} & \textcolor{white}{\textbf{Punca}} & \textcolor{white}{\textbf{Penyelesaian}} \\
\midrule
\endhead
\rowcolor{tablebg}
Tidak terima notifikasi & Telegram tidak dihubungkan & 1) Hubungkan Telegram 2) Semak spam \\
Pautan ``Sudah Tamat Tempoh'' & Pautan telah digunakan & Log masuk, jawab via ``Tugasan Lantikan'' \\
\rowcolor{tablebg}
``Tugasan Sudah Dijawab'' & Sudah jawab saluran lain & Tiada tindakan perlu \\
Ingin tukar jawapan & Jawapan tidak boleh diubah & Hubungi Admin \\
\bottomrule
\end{longtable}

\section{Masalah Penilaian}

\begin{longtable}{L{4cm} L{3.5cm} L{6.5cm}}
\toprule
\rowcolor{tableheader}
\textcolor{white}{\textbf{Masalah}} & \textcolor{white}{\textbf{Punca}} & \textcolor{white}{\textbf{Penyelesaian}} \\
\midrule
\endhead
\rowcolor{tablebg}
Tiada perlawanan ``Perlu Dinilai'' & Belum terima/jawab lantikan & Semak ``Tugasan'', tekan ``Terima'' \\
Tidak boleh hantar laporan & Markah belum lengkap & Isi markah semua pegawai \\
\rowcolor{tablebg}
``Laporan Sudah Dihantar'' & Dihantar kaedah lain & Hubungi Admin \\
Borang token tidak boleh dibuka & Token tidak sah/tamat & Hubungi Admin untuk pautan baharu \\
\bottomrule
\end{longtable}

\section{Masalah Profil \& Telegram}

\begin{longtable}{L{4cm} L{3.5cm} L{6.5cm}}
\toprule
\rowcolor{tableheader}
\textcolor{white}{\textbf{Masalah}} & \textcolor{white}{\textbf{Punca}} & \textcolor{white}{\textbf{Penyelesaian}} \\
\midrule
\endhead
\rowcolor{tablebg}
Gambar profil tidak bertukar & Cache pelayar & Muat semula (Ctrl+Shift+R) \\
``Gagal memuat naik gambar'' & Format/saiz & JPEG/PNG/WebP, maks 5MB \\
\rowcolor{tablebg}
Bot Telegram tidak bertindak balas & Pautan tamat & Profil $\rightarrow$ ``Hubungkan Telegram'' semula \\
Status ``Tidak Dihubungkan'' & Proses gagal & Cuba klik pautan sekali lagi \\
\rowcolor{tablebg}
Tidak terima mesej Telegram & Bot disekat & Buka semula perbualan, pastikan tidak disekat \\
\bottomrule
\end{longtable}

\section{Masalah Umum}

\begin{longtable}{L{5cm} L{9cm}}
\toprule
\rowcolor{tableheader}
\textcolor{white}{\textbf{Masalah}} & \textcolor{white}{\textbf{Penyelesaian}} \\
\midrule
\endhead
\rowcolor{tablebg}
Halaman tidak dimuatkan & 1) Muat semula 2) Kosongkan cache 3) Cuba pelayar lain \\
Data tidak dikemaskini & Muat semula halaman \\
\rowcolor{tablebg}
Butang tidak berfungsi & 1) Pastikan JavaScript aktif 2) Kosongkan cache \\
Ralat ``500'' / ``Ralat Pelayan'' & Cuba lagi selepas beberapa minit. Hubungi Admin jika berterusan \\
\rowcolor{tablebg}
Paparan tidak kemas (telefon) & Kemaskini pelayar. Cuba mod landskap untuk jadual \\
\bottomrule
\end{longtable}

\vspace{1cm}

\begin{center}
\textcolor{pbnpgold}{\rule{8cm}{2pt}}
\end{center}

\vspace{0.5cm}

\section*{Sokongan Teknikal}

Jika masalah tidak dapat diselesaikan, sila hubungi:

\begin{table}[H]
\centering
\begin{tabularx}{0.7\textwidth}{l X}
\toprule
\rowcolor{tableheader}
\textcolor{white}{\textbf{Saluran}} & \textcolor{white}{\textbf{Maklumat}} \\
\midrule
\rowcolor{tablebg}
Pentadbir Sistem & Hubungi melalui sistem (notifikasi) \\
Emel Sokongan & Hubungi Admin melalui emel rasmi PBNP \\
\bottomrule
\end{tabularx}
\end{table}

\vspace{2cm}

\begin{center}
{\small\textcolor{pbnpblue!50}{Dokumen ini dijana dari analisis sistem RefPahang versi 3.0}}\\
{\small\textcolor{pbnpblue!50}{Tarikh: 5 April 2026}}\\[0.3cm]
{\small\textcolor{pbnpblue!50}{\textcopyright\ Persatuan Bola Sepak Negeri Pahang (PBNP)}}
\end{center}

\end{document}
